% Introduction Section for MobilePIXORNeXt
% Written according to Springer Nature - Arabian Journal for Science and Engineering (AJSE)
% Structure: 5-Paragraph Funnel (Context & Significance -> Existing Approaches & Bottlenecks -> Motivation & Key Insight -> Proposed Method Overview -> Main Contributions)

Three-dimensional (3D) spatial perception from Light Detection and Ranging (LiDAR) point clouds provides essential metric localization for autonomous navigation, mobile robotics, and intelligent transportation systems. Unlike passive optical cameras that suffer from projective depth ambiguity and ambient lighting fluctuations, LiDAR sensors emit laser pulses to capture metric geometric surfaces during day and night~\cite{geiger2012kitti}. Autonomous vehicle perception requires detecting, classifying, and orienting surrounding objects, such as passenger vehicles, pedestrians, and cyclists, within stringent latency constraints. In operational systems, end-to-end execution must generally remain below $100\,\text{ms}$ per LiDAR scan~\cite{yang2018pixor}. Deploying accurate 3D detectors on onboard edge platforms requires balancing spatial localization precision against computing hardware budgets. Embedded accelerators have restricted memory bandwidth, cache sizes, and arithmetic capacity under strict thermal envelopes. Compact LiDAR detectors must sustain high throughput while preserving geometric precision.

Existing 3D LiDAR detectors present distinct trade-offs between geometric precision and edge computational feasibility across several architectural paradigms. Volumetric voxel networks such as Sparsely Embedded Convolutional Detection (SECOND)~\cite{yan2018second} and hybrid two-stage frameworks (e.g., PointRCNN~\cite{shi2019pointrcnn}, Part-$A^2$~\cite{shi2020part}, and PV-RCNN~\cite{shi2020pvrcnn}) rely on 3D sparse convolutions or continuous point abstractions. While achieving high accuracy, their large parameter counts ($4.0\text{--}63.8\,\text{M}$) and substantial full-pipeline latencies ($93.7\text{--}194.5\,\text{ms}$) restrict deployment on resource-constrained platforms. Point-based direct architectures such as PointNet++~\cite{qi2017pointnetplusplus} process raw coordinates directly without discretization, but iterative Farthest Point Sampling and ball query neighbor searches impose $O(N^2)$ memory bandwidth demands that saturate edge caches. In contrast, Range-View spherical projection methods such as LaserNet~\cite{meyer2019lasernet} project 3D point clouds onto cylindrical range images to apply standard 2D convolutions; however, spherical projection causes severe perspective scale dilation with distance and boundary blurring across foreground-background occlusions. Pillar-based detectors, including PointPillars~\cite{lang2019pointpillars}, PillarNeXt~\cite{li2023pillarnext}, and PillarNeSt~\cite{mao2023pillarnest}, project points into vertical columns and utilize PointNet encoders to construct 2D Bird's-Eye View (BEV) representations. They achieve higher execution speed than volumetric networks, but point gathering, dynamic memory allocation, and sparse scattering introduce runtime and PCIe memory transfer overhead. Proposal-free BEV architectures such as PIXOR~\cite{yang2018pixor} and MobilePIXOR~\cite{nguyen2024indoor} replace learned encoders with deterministic vertical occupancy grids. However, binary occupancy grids on uniform $0.1\,\text{m}$ grids discard continuous subcell centroid displacements and vertical elevation dispersion moments. This spatial quantization causes boundary staircasing and localization jitter, particularly for small objects. Recent extensions such as PillarHist~\cite{zhou2025pillarhist} and TriBand-BEV~\cite{khoshkdahan2026triband} explore vertical histograms or multi-band intensity slicing, but existing compact representations generally omit continuous subcell coordinate offsets alongside higher-order vertical spatial moments. Additional bottlenecks include synthetic placement artifacts from naive cut-and-paste data augmentation, restricted receptive fields in small convolutional backbones, and uncalibrated classification heatmaps that allow poorly oriented candidates to suppress well-localized boxes during Non-Maximum Suppression (NMS)~\cite{li2020gfl}. Compact perception systems must resolve these representation, augmentation, receptive field, and ranking challenges within a sub-million-parameter budget.

Vehicular collision avoidance and kinematic trajectory planning operate predominantly on the Special Euclidean manifold $SE(2)$ (planar translation and yaw heading). In on-road autonomous driving, horizontal clearance margins and Time-to-Collision (TTC) along the ground plane dictate the vast majority of critical control decisions. Therefore, predicting high-precision oriented 2D bounding footprints on a metric Bird's-Eye View (BEV) plane provides the most computationally direct formulation for edge vehicular safety. We investigate whether a deterministic, parameter-free BEV representation can retain fine sub-voxel geometry without learned point encoders. By combining multi-stratum vertical densities, continuous elevation moments, surface reflectivity statistics, and continuous horizontal centroid offsets, a deterministic tensor preserves spatial boundaries without encoder training overhead. Our core research question is: \textit{Can an ultra-compact LiDAR detector operating on a parameter-free BEV tensor approach the localization accuracy of heavy volumetric and pillar architectures while achieving higher end-to-end throughput on edge hardware?} To address this question, we recognize that compact backbones cannot succeed through isolated architectural scaling alone; rather, they require an end-to-end co-design between sensor physical modeling and geometric loss supervision. We pair this representation with a single-branch ConvNeXt-style backbone~\cite{liu2022convnext} employing $7 \times 7$ depthwise convolutions and parameter-free 3D spatial-energy attention (SimAM)~\cite{yang2021simam}, expanding physical receptive fields to cover complete passenger vehicles at stride 8 without multi-branch training overhead. Physics-consistent synthesis (PC-GTA) regularizes distant returns by enforcing ground elevation flushness and raycast shadows, while multi-axis geometric supervision (OGA) and localization-quality ranking (IQA) align training losses and detection confidence with rotated bounding box overlap. These design choices collectively narrow the performance gap between compact 2D BEV representations and multi-million-parameter 3D sparse networks.

We propose \textbf{MobilePIXORNeXt}, a compact, proposal-free, multi-class BEV LiDAR object detector with $753{,}569$ trainable detector parameters ($<0.76\,\text{M}$). Figure~\ref{fig:r37_overview} illustrates the end-to-end pipeline. The parameter-free \textbf{Strata-BEV14} rasterizer converts an unordered point cloud into a 14-channel statistical BEV feature map ($800 \times 704$ at $0.1\,\text{m}$ grid resolution) using CPU execution. Point-cloud loading and rasterization take $20.94\,\text{ms}$ in the profiled pipeline. The 14 channels capture vertical stratum densities, continuous elevation moments ($z_{\max}, z_{\mathrm{mean}}, z_{\mathrm{span}}, z_{\mathrm{std}}$), reflectance statistics, and subcell centroid displacement offsets $(\Delta\bar{x}, \Delta\bar{y})$ to mitigate quantization errors. \textbf{Physics-Consistent Ground-Truth Augmentation (PC-GTA)} enforces local road elevation fitting, static obstacle clearance, and range-adaptive point downsampling to prevent unphysical synthesis artifacts. The \textbf{MobilePIXORNeXt-S} backbone extracts multi-scale features across stages $C_2, C_3, C_4$ using $7 \times 7$ depthwise inverted bottleneck blocks ($e=2.5$), SiLU activations, LayerScale stabilization, and parameter-free \textbf{Block-SimAM} energy attention embedded within all nine residual blocks. The top-down \textbf{Scale-Gated Feature Pyramid Network (SG-FPN Neck, $F_{32}$)} balances semantic abstraction and fine spatial details via bilinear upsampling, depthwise refinement at stride 8, standard projection at stride 4, and learnable depthwise scale gates, avoiding transposed-convolution checkerboard patterns. Car ($K_1=1$) and Pedestrian/Cyclist ($K_2=2$) predictions are decoupled into independent \textbf{GroupedHeader} towers sharing an initial $3 \times 3$ Conv-BN-SiLU block. Training utilizes \textbf{Oriented Geometric Alignment (OGA)}, which combines Multi-Axis Projection Generalized IoU (MGIoU)~\cite{le2026mgiou} with scale-normalized $\pi$-symmetric corner distance, balanced across tasks by 12-parameter \textbf{Temperature-Softmax Bounded Uncertainty Weighting (T-SBUW)}. At inference, \textbf{IoU-Aware Quality Alignment (IQA)} evaluates composite detection scores $s_{\mathrm{det}} = (p_{\mathrm{cls}})^{1-\alpha} \hat{q}^\alpha$ prior to class-wise rotated Non-Maximum Suppression (NMS). For memory-critical platforms, we also present \textbf{Q-OGA}, a head-free quality-coupled objective that eliminates the separate IQA prediction tower.

The principal contributions of this work are summarized as follows:
\begin{itemize}
    \item \textbf{Ultra-compact proposal-free BEV detection architecture.} We design \textbf{MobilePIXORNeXt}, a compact ($753{,}569$ trainable parameters) LiDAR detector that eliminates learned PointNet encoders and 3D sparse convolutions. The architecture integrates the 14-channel deterministic \textbf{Strata-BEV14} representation with the \textbf{MobilePIXORNeXt-S} backbone and parameter-free \textbf{Block-SimAM} energy attention. The \textbf{SG-FPN} neck ($F_{32}$) fuses multi-scale features for decoupled \textbf{GroupedHeader} towers, resolving scale competition across road object classes.
    
    \item \textbf{Physics-consistent augmentation, geometric supervision, and head-free alternative.} We develop \textbf{PC-GTA} to enforce terrain alignment, obstacle clearance, and distance-dependent point subsampling during training. We formulate \textbf{OGA}, coupling Multi-Axis Projection GIoU with scale-normalized $\pi$-symmetric corner supervision, balanced across task objectives via 12-parameter \textbf{T-SBUW}. Confidence scores are calibrated via \textbf{IQA} quality alignment. In addition, we develop \textbf{Q-OGA} as an ultra-compact, head-free alternative that integrates localization quality directly into classification targets, saving $18{,}626$ parameters while attaining $69.20\%$ Moderate mAP.
    
    \item \textbf{Empirical validation and efficiency benchmarking.} Evaluated on the KITTI validation benchmark under official 40-point interpolated BEV AP ($R_{40}$), MobilePIXORNeXt establishes a competitive accuracy, model footprint, and runtime trade-off. It achieves \textbf{71.46\% Moderate mAP} and \textbf{73.42\% $\text{mAP}_9$} with \textbf{0.754\,M parameters} ($9.34\,\text{MB}$ checkpoint). It outperforms PointPillars ($70.13\%$) by \textbf{1.33 percentage points} with an \textbf{84.41\% parameter reduction}, and trails SECOND ($72.09\%$) by only \textbf{0.63 percentage points} while utilizing \textbf{85.85\% fewer parameters} ($7.07\times$ smaller parameter footprint). On Pedestrian Moderate AP, it reaches \textbf{60.75\%} (+4.74\,pp over PointPillars, +4.19\,pp over SECOND). The complete end-to-end pipeline executes in \textbf{60.96\,ms} (\textbf{16.41\,FPS}) on an NVIDIA Tesla T4 GPU, delivering the highest end-to-end throughput among the evaluated real-time detectors. In ablation studies, PC-GTA improves accuracy by \textbf{+6.03\,pp} ($60.22\%$ to $66.25\%$), OGA and IQA add \textbf{+3.20\,pp}, and Block-SimAM adds \textbf{+0.88\,pp} with zero additional parameters.
\end{itemize}

Section~\ref{sec:related_work} reviews LiDAR representations, lightweight backbones, and geometric losses. Section~\ref{sec:method} details the MobilePIXORNeXt architecture and loss formulations. Section~\ref{sec:experiments} specifies experimental settings. Section~\ref{sec:results} reports benchmark comparisons, latency breakdowns, ablation studies, and qualitative results. Section~\ref{sec:conclusion} concludes the paper.
